% cp_t2_panelB — Replication of Cieslak–Povala (2015), Table 2, Panel B:
% predictive R2 under the expectations-hypothesis null.
% Numbers transcribed from tables/cp_t2_panelB.pdf (generated by empirical.R).
\small
\begin{tabular}{cc ccc}
  \toprule
  & & \multicolumn{3}{c}{Percentiles of $\bar{R}^{2}$} \\
  \cmidrule(lr){3-5}
  $\phi_{\tau}$ & $\phi_{r}$ & 5th & 50th & 95th \\
  \midrule
  $0.800$ & $0.75$ & $-0.003$ & $0.007$ & $0.045$ \\
  $0.975$ & $0.75$ & $-0.002$ & $0.019$ & $0.096$ \\
  $0.999$ & $0.75$ & $-0.002$ & $0.022$ & $0.102$ \\
  \midrule
  $0.975$ & $0.60$ & $-0.003$ & $0.015$ & $0.083$ \\
  $0.975$ & $0.75$ & $-0.002$ & $0.019$ & $0.102$ \\
  $0.975$ & $0.90$ & $-0.002$ & $0.029$ & $0.125$ \\
  \bottomrule
\end{tabular}
\tabnotes{Each row reports the $5$th, $50$th, and $95$th percentiles of the
  adjusted $R^{2}$ from the predictive regression of $\rxbar_{t+12}$ on the
  latent factors $(\tau_{t},r_{t})$ in data simulated from the calibrated
  expectations-hypothesis economy of \Cref{sec:meth-ehmc}, in which risk premia
  are constant ($\lambda=0$); $5{,}000$ simulations of length $T=470$ months,
  the sample length of \citet{cieslak2015}. The upper block varies the
  persistence of trend inflation $\phi_{\tau}$ holding $\phi_{r}=0.75$; the
  lower block varies the persistence of the real-rate factor $\phi_{r}$
  holding $\phi_{\tau}=0.975$ (the blocks are simulated independently, so the
  shared grid point differs slightly between them). The $95$th percentile
  rises with both persistences, reproducing the shape of the null distribution
  in \citet{cieslak2015}, Table~2, Panel~B; its level lies below their
  $0.19$--$0.23$ because the original calibration is underspecified, a gap
  documented in \Cref{sec:meth-ehmc} rather than closed.}
